% ============================================================
%  Paper B -- draft manuscript
%  Target: Medical Image Analysis (swap \documentclass to
%  elsarticle when submitting) or TMLR.
%
%  PLACEHOLDERS marked [PENDING] must be filled before
%  submission. Search the file for "PENDING".
% ============================================================
\documentclass[11pt,a4paper]{article}

\usepackage[utf8]{inputenc}
\usepackage[T1]{fontenc}
\usepackage{amsmath,amssymb}
\usepackage{booktabs}
\usepackage{graphicx}
\usepackage{xcolor}
\usepackage[margin=1in]{geometry}
\usepackage{caption}
\usepackage{subcaption}
\usepackage[hidelinks]{hyperref}
\usepackage{siunitx}
\sisetup{detect-weight=true}

\newcommand{\pending}[1]{\textcolor{red}{[PENDING: #1]}}
\newcommand{\lam}{\lambda_{\mathrm{adv}}}

% ============================================================
%  FILL-IN BLOCK -- every outstanding value lives here.
%  Replace each \pending{...} with the number; the text and
%  tables update themselves. Run scripts/emit_tex_values.py on
%  the HPC to print these lines ready to paste.
% ============================================================
\newcommand{\xdomHalf}{0.270\,$\pm$\,0.034}
\newcommand{\xdomOne}{0.261\,$\pm$\,0.016}
\newcommand{\xdomFour}{0.276\,$\pm$\,0.025}
\newcommand{\xdomEight}{0.269\,$\pm$\,0.023}
\newcommand{\xdomEffb}{0.298\,$\pm$\,0.010}
\newcommand{\nAK}{202}      % AK count in pad\_heldout
\newcommand{\nSCC}{66}     % SCC count in pad\_heldout
\newcommand{\nMEL}{9}               % melanoma count in pad\_heldout
\newcommand{\repolink}{\pending{repository link}}
% ============================================================

% Degrade gracefully when a figure file is not present, so the manuscript
% compiles before the figures are copied in. Replace with plain
% \includegraphics once results/paperB/figures/ is available locally.
\newcommand{\figorbox}[2]{%
  \IfFileExists{#2}{\includegraphics[width=#1]{#2}}{%
    \fbox{\parbox[c][45mm][c]{0.96\linewidth}{\centering\ttfamily
    \small missing figure:\\\detokenize{#2}}}}}

\title{\textbf{Confidently Wrong on Unseen Domains:\\
Domain-Adversarial Training Silently Disables\\
Out-of-Distribution Monitoring}}

\author{\pending{author list and affiliations}}
\date{}

\begin{document}
\maketitle

% ------------------------------------------------------------
\begin{abstract}
\noindent\textbf{Background.} Domain-adversarial invariance training is widely
recommended for removing acquisition shortcuts from medical imaging models.
Its effect on the out-of-distribution (OOD) monitoring that such models rely
on at deployment has not been measured.

\noindent\textbf{Methods.} We trained a dual-encoder factorization model on
ISIC~2019 (dermoscopy, 25{,}331 images) and PAD-UFES-20 (clinical photographs,
2{,}298 images), sweeping the adversarial weight $\lam$ across seven values
with three to five seeds each (27 runs). At every $\lam$ we measured domain
leakage by linear probe, in-distribution balanced accuracy, calibration on
in-distribution and out-of-domain data, OOD detection with five scoring rules,
and zero-shot cross-domain diagnostic accuracy on a patient-level held-out
PAD split. We validated on Fitzpatrick17k, a third dataset never seen by the
adversary.

\noindent\textbf{Results.} Domain leakage fell to its chance floor at the
smallest non-zero weight tested (balanced accuracy $0.915 \to 0.553$, floor
$0.5$), while in-distribution balanced accuracy did not degrade
($0.692 \to 0.707$). Over the same range, OOD detection fell from
AUROC~$0.863$ to $0.508$ and inverted below chance ($0.413$), with OOD samples
ranked as more in-distribution than in-distribution data throughout the score
distribution rather than only in the tails. Cross-domain diagnostic accuracy
did not improve at any $\lam$ (balanced accuracy $0.249$--$0.291$ against a
six-class floor of $0.167$, and below the $0.298$ of a non-adversarial control
at every setting), and predictions collapsed toward the benign majority class:
actinic keratosis recall fell from $0.244$ to $0.016$ ($n = 202$). Mean confidence on
never-seen-domain images rose from $0.489$ to $0.930$, overtaking
in-distribution confidence, so that out-of-domain calibration error tripled
($\mathrm{ECE}\ 0.247 \to 0.746$) while in-distribution ECE stayed flat at
$\approx 0.10$. Leakage did not mediate the OOD collapse: it saturated at
chance while AUROC kept falling, and at the backbone leakage was flat while
AUROC fell $0.749 \to 0.475$. Six candidate explanations were tested and
refuted.

\noindent\textbf{Conclusions.} When domain-adversarial training succeeds,
models become confidently wrong on unseen domains and their OOD monitoring
fails, while in-distribution metrics improve. OOD detection and out-of-domain
calibration should be reported as a function of achieved invariance rather
than at a single operating point.
\end{abstract}

\noindent\textbf{Keywords:} domain-adversarial training; out-of-distribution
detection; calibration; distribution shift; skin lesion classification;
model monitoring

% ------------------------------------------------------------
\section{Introduction}

Deep models for skin lesion classification are known to exploit acquisition
artefacts rather than pathology: ruler marks, ink annotations, imaging
hardware and site identity all correlate with diagnosis in public benchmarks
\cite{winkler2019,bissoto2020,daneshjou2022}. The standard remedy is to make
the representation invariant to the acquisition domain, most often through
domain-adversarial training with a gradient reversal layer
\cite{ganin2016,bousmalis2016}. The stated purpose of this remedy is safe
deployment across sites.

Deployment across sites also relies on a second mechanism: out-of-distribution
monitoring. A deployed system is expected to flag inputs unlike its training
data so that a human reviews them rather than the model acting alone. This
paper asks what invariance training does to that mechanism.

The answer is that it destroys it, and does so invisibly. Sweeping the
adversarial weight in a controlled setting, we find that OOD detection
collapses from AUROC $0.863$ to $0.508$ at the smallest non-zero weight tested
and then inverts below chance, so that images from an unseen clinic are ranked
as \emph{more} routine than familiar ones. Over the same range the model
becomes more confident on those unseen images than on its own held-out test
data, and its predictions collapse toward the benign majority class. None of
this is visible to the metrics a developer monitors: in-distribution accuracy
does not fall, in-distribution calibration does not change, and the domain
leakage probe --- the very metric the technique is adopted to improve ---
improves as intended.

The irony is the point. The technique is adopted for cross-site robustness and
it degrades precisely the cross-site safety machinery, while every monitored
number moves in the direction its user is hoping for.

\paragraph{Contributions.}
\begin{enumerate}
\item A controlled intervention: a seven-point sweep of the adversarial weight
with 27 training runs in a fixed architecture, isolating the effect of
invariance pressure from architectural and data confounds.
\item The safety failure: OOD detection collapses as a cliff rather than a
trade-off curve, inverts below chance, and does so on a third domain the
adversary never saw.
\item The invisibility: in-distribution accuracy, in-distribution calibration
and the domain leakage probe all fail to reveal it.
\item Evidence that the leakage probe does not mediate the collapse, so it
cannot serve as a proxy for it.
\item Six candidate explanations tested and refuted, plus a pre-registered
synthetic reproduction that failed --- delimiting the phenomenon for future
work.
\item A reporting protocol that makes the failure visible at near-zero cost.
\end{enumerate}

% ------------------------------------------------------------
\section{Related work}

\paragraph{Shortcut learning in dermatology.}
Winkler et al.~\cite{winkler2019} showed that ink markings shift melanoma
predictions; Bissoto et al.~\cite{bissoto2020} showed that classifiers retain
competitive accuracy on ISIC data when the lesion itself is masked;
Daneshjou et al.~\cite{daneshjou2022} documented degraded performance on
darker skin tones. These motivate invariance-based remedies.

\paragraph{Invariant representation learning.}
Domain-adversarial training \cite{ganin2016} and domain separation networks
\cite{bousmalis2016} remove domain information by adversarial or
orthogonality constraints. \textbf{The architecture used here is
equivalent to a domain separation network and is not a contribution of this
paper.} It is a vehicle for the intervention.

\paragraph{Covariate shift versus semantic novelty.}
Full-spectrum OOD detection \cite{yang2023} argues that the OOD literature
conflates two distinct shifts, and proposes a four-way taxonomy: training ID,
\emph{covariate-shifted ID}, near-OOD and far-OOD. Crucially, it prescribes
that covariate-shifted in-distribution samples should be \textbf{tolerated as
in-distribution rather than flagged}, on the grounds that rejecting them
contradicts the goal of generalization. OpenOOD v1.5 \cite{zhang2023} adopts
the same full-spectrum setting at ImageNet scale and for foundation models.

Our setting falls squarely in the covariate-shifted-ID cell: PAD-UFES-20
contains the same diagnostic classes as ISIC, photographed differently. A
reader of \cite{yang2023} might therefore regard a model that accepts these
images as in-distribution as behaving \emph{correctly}.

We argue this is exactly the gap. \textbf{Tolerance of covariate shift is safe
only when accompanied by generalization, and the two are evaluated
separately or not at all.} Our models achieve the prescribed tolerance --- by
$\lam = 2$ they rank out-of-domain images as more in-distribution than their
own test data --- while cross-domain balanced accuracy sits at $0.249$ against
a floor of $0.167$, plain accuracy is below the majority-class rate, and
confidence on those images exceeds in-distribution confidence. Tolerating
inputs one cannot classify is worse than flagging them, and no current
protocol separates the two cases. We therefore report acceptance and
cross-domain accuracy jointly throughout.

\paragraph{Calibration as a route to invariance --- and the reverse.}
Ovadia et al.~\cite{ovadia2019} showed that calibration degrades under shift
for ordinary training. Wald et al.~\cite{wald2021} went further and
established the forward direction: enforcing \emph{multi-domain calibration}
yields invariance, with models so calibrated provably free of certain spurious
correlations, and they propose a calibration loss (CLOvE) validated on WILDS.

Our results show the reverse implication does not hold. Obtaining invariance
adversarially does not yield calibration; it \emph{destroys} out-of-domain
calibration in a dose-dependent way while leaving in-distribution calibration
untouched. Calibration and adversarial invariance are therefore not
interchangeable routes to the same representation, and the direction of
enforcement matters.

\paragraph{What is new here.}
(i) The \emph{shape} of the dependence --- a cliff at the smallest effective
adversarial weight rather than a graded trade-off; (ii) \emph{sub-chance}
inversion, in which out-of-domain data is ranked as more typical than
in-distribution data throughout the score distribution; (iii) the
demonstration that the failure is invisible to every routinely monitored
metric; (iv) the finding that the standard domain leakage probe neither
mediates nor predicts it; (v) the consequence for calibration and for
diagnostic decisions on out-of-domain data; and (vi) the observation that
full-spectrum tolerance is achieved here without the generalization that would
make it safe.

% ------------------------------------------------------------
\section{Materials and methods}

\subsection{Data}

\textbf{ISIC 2019} \cite{isic2019}: 25{,}331 dermoscopic images, eight
diagnostic classes (MEL, NV, BCC, AK, BKL, DF, VASC, SCC), CC~BY-NC~4.0.
\textbf{PAD-UFES-20} \cite{pacheco2020}: 2{,}298 smartphone clinical
photographs, six classes mapping into the same label space (DF and VASC
absent), CC~BY~4.0. \textbf{Fitzpatrick17k} \cite{groh2021}: 3{,}887 of
16{,}577 images obtained; annotations CC~BY-NC-SA~3.0, images web-scraped
under mixed copyright, used for evaluation only and not redistributed.

Images were cropped to a soft lesion-centred bounding box via Otsu
thresholding on the deviation from the border median, expanded by a 0.30
margin with a 0.70 minimum crop ratio, and resized to $224\times224$.

\subsection{Splits and the PAD held-out set}

ISIC was split $60/20/20$ into train/validation/test, stratified by label
(\texttt{random\_state}~$=42$).

\textbf{PAD appears in training and would otherwise also be the OOD set.}
The classification loss never sees PAD labels (\texttt{ignore\_index}$=-1$),
but PAD images enter the context and adversarial branches. To remove the
resulting circularity we partitioned PAD \emph{at patient level} into
\texttt{pad\_adv} (1{,}582 images, 961 patients), visible to those branches
during training, and \texttt{pad\_heldout} (716 images, 412 patients), never
seen in any form. No patient and no lesion appears in both. \textbf{All OOD
and cross-domain numbers reported in the main text use
\texttt{pad\_heldout}.} Fitzpatrick17k provides a second never-adversarial
domain.

\subsection{Architecture and objective}

A dual-encoder model with EfficientNet-B3 backbones produces a lesion latent
$z_{\ell} \in \mathbb{R}^{16}$ (batch-normalised, denoted
$z_{\ell}^{\mathrm{norm}}$) and a context latent $z_{c} \in \mathbb{R}^{64}$.
Diagnosis uses $z_{\ell}^{\mathrm{norm}}$ only, through a linear head. The
objective is
\begin{equation}
\mathcal{L} = \mathcal{L}_{\mathrm{cls}}
+ \lambda_{\mathrm{ctx}}\mathcal{L}_{\mathrm{ctx}}
+ \lam\,\mathcal{L}_{\mathrm{adv}}
+ \lambda_{\mathrm{orth}}\mathcal{L}_{\mathrm{orth}},
\end{equation}
with $\mathcal{L}_{\mathrm{cls}}$ on ISIC only, $\mathcal{L}_{\mathrm{ctx}}$ a
domain cross-entropy on $z_c$, $\mathcal{L}_{\mathrm{adv}}$ a domain
cross-entropy on $z_{\ell}^{\mathrm{norm}}$ behind a gradient reversal layer,
and $\mathcal{L}_{\mathrm{orth}}$ a squared-cosine penalty between the two
latents ($\lambda_{\mathrm{orth}} = 1$ throughout).

AdamW, base learning rate $10^{-4}$, adversary learning rate multiplier $30$,
classifier multiplier $0.2$, 40 epochs, batches balanced across domains,
checkpoint selected on validation accuracy.

\subsection{The intervention}

$\lam \in \{0,\,0.25,\,0.5,\,1,\,2,\,4,\,8\}$; five seeds at
$\lam \in \{0, 2, 8\}$ and three at the interior points, for 27 runs.
Everything else is held fixed, so $\lam$ is the only manipulated variable.

\subsection{Measurements}

\paragraph{Domain leakage.} Logistic probe on frozen features predicting the
source dataset, reported as \textbf{balanced accuracy, whose chance floor is
$0.5$ for this two-domain problem}. Earlier work on this dataset pair reports
plain accuracy against a majority floor of $0.688$; the two scales must not be
compared to each other. Both are stated where relevant.

\paragraph{OOD detection.} MSP, free energy, maximum cosine to a class
centroid, class-conditional Mahalanobis with shared covariance, and $k$-NN
distance ($k \in \{10,50,200\}$). All statistics are fit on the ISIC training
split alone. \textbf{These five rules are not independent evidence}: the
classifier is a linear map from $z_{\ell}^{\mathrm{norm}}$, so all five are
functions of the same 16-dimensional vector. Independent evidence comes from
repeating the analysis on the pre-projection backbone representation
(Section~\ref{sec:nonmediation}).

\paragraph{Calibration.} Expected calibration error, 15 bins, computed
separately on ISIC test and on \texttt{pad\_heldout} restricted to the six
shared classes. \textbf{Note that out-of-domain ECE here is arithmetically
close to mean confidence minus accuracy}, so the confidence, accuracy and ECE
results in Section~\ref{sec:confidently} are three views of one phenomenon,
not three independent findings.

\paragraph{Cross-domain diagnosis.} Zero-shot evaluation on
\texttt{pad\_heldout} over the six shared classes. PAD labels were never used
for $\mathcal{L}_{\mathrm{cls}}$.

% ------------------------------------------------------------
\section{Results}

\subsection{The dial}

Table~\ref{tab:dial} reports every measurement as a function of $\lam$.
Figure~\ref{fig:dial} shows the four quantities that matter most on a shared
axis.

\begin{table}[t]
\centering
\small
\caption{The intervention. Mean $\pm$ s.d.\ across seeds. Leakage is balanced
accuracy (chance $0.5$); cross-domain accuracy is balanced over six classes
(chance $0.167$). Arrows mark the direction a practitioner would call
favourable.}
\label{tab:dial}
\resizebox{\textwidth}{!}{%
\begin{tabular}{lcccccccc}
\toprule
$\lam$ & $n$ & Leakage $\downarrow$ & ID bal.\ acc $\uparrow$ & ID ECE
& OOD ECE & OOD AUROC & Fitz.\ AUROC & X-dom bal.\ acc \\
\midrule
0    & 5 & 0.915\,$\pm$\,0.021 & 0.692 & 0.107 & 0.247 & \textbf{0.863\,$\pm$\,0.025} & 0.764\,$\pm$\,0.038 & 0.291\,$\pm$\,0.016 \\
0.25 & 3 & 0.553\,$\pm$\,0.008 & 0.696 & 0.101 & 0.684 & \textbf{0.508\,$\pm$\,0.022} & 0.439\,$\pm$\,0.061 & 0.291\,$\pm$\,0.020 \\
0.5  & 3 & 0.555\,$\pm$\,0.029 & 0.701 & 0.096 & 0.724 & 0.449\,$\pm$\,0.008 & 0.449\,$\pm$\,0.025 & \xdomHalf \\
1    & 3 & 0.569\,$\pm$\,0.010 & 0.701 & 0.097 & 0.733 & 0.439\,$\pm$\,0.040 & 0.457\,$\pm$\,0.088 & \xdomOne \\
2    & 5 & 0.584\,$\pm$\,0.027 & 0.707 & 0.099 & 0.746 & \textbf{0.413\,$\pm$\,0.021} & 0.391\,$\pm$\,0.068 & 0.249\,$\pm$\,0.030 \\
4    & 3 & 0.548\,$\pm$\,0.035 & 0.686 & 0.100 & 0.726 & 0.443\,$\pm$\,0.038 & 0.431\,$\pm$\,0.036 & \xdomFour \\
8    & 5 & 0.625\,$\pm$\,0.045 & 0.679 & 0.099 & 0.713 & 0.481\,$\pm$\,0.026 & 0.474\,$\pm$\,0.060 & \xdomEight \\
\midrule
\multicolumn{2}{l}{EffNet-B3 control} & 0.802 & 0.683 & 0.096 & --- & 0.726\,$\pm$\,0.021 & 0.649 & \xdomEffb \\
\multicolumn{2}{l}{ResNet-50 baseline} & 0.979 & 0.655 & 0.101 & --- & 0.868\,$\pm$\,0.019 & 0.994 & 0.391 \\
\bottomrule
\end{tabular}}
\end{table}

\begin{figure}[t]
\centering
\figorbox{\textwidth}{results/paperB/figures/fig1_dial_four_panel.pdf}
\caption{\textbf{The intervention, four views, one axis.} Leakage falls to its
chance floor and in-distribution balanced accuracy does not degrade --- both
favourable. Over the same range OOD AUROC collapses through chance and
inverts, and out-of-domain calibration error triples. A developer monitoring
the first two panels would see nothing wrong.}
\label{fig:dial}
\end{figure}

\subsection{Invariance is achieved, and in-distribution metrics do not warn}

Leakage falls from $0.915$ to $0.553$ at $\lam = 0.25$, against a chance floor
of $0.5$ --- that is, essentially complete invariance by the smallest non-zero
weight tested. In-distribution balanced accuracy does not degrade over the
useful range ($0.692 \to 0.707$; we claim no improvement, the difference being
small relative to seed variance) and in-distribution ECE is flat at
$\approx 0.10$ at every $\lam$.

\subsection{OOD detection collapses as a cliff, and inverts}

Mahalanobis AUROC falls $0.863 \to 0.508$ between $\lam = 0$ and
$\lam = 0.25$, and reaches $0.413$ at $\lam = 2$. \textbf{The collapse is
essentially complete at the smallest non-zero weight}; between $\lam = 0.5$
and $\lam = 8$ the value varies only between $0.413$ and $0.481$. This is a
cliff, not a trade-off curve, and there is no safe operating point on it.

All five scoring rules move together (MSP $0.970 \to 0.437$, energy
$0.983 \to 0.441$, cosine $0.901 \to 0.435$, $k$-NN $0.941 \to 0.424$, stable
across $k \in \{10,50,200\}$). As noted in the methods, these are five views of
one 16-dimensional representation and are not independent.

The inversion is a \textbf{bulk location shift, not a tail artefact}. At
$\lam = 2$, removing the pooled $90$th, $95$th or $99$th percentile leaves
AUROC at $0.412$, $0.416$ and $0.424$ against $0.427$ for the full set, and the
medians themselves invert (ISIC $13.8$ versus \texttt{pad\_heldout} $10.0$).
Out-of-domain images are scored as more typical than in-distribution images
throughout the distribution.

It is not a numerical artefact of the covariance estimate: Ledoit--Wolf
shrinkage leaves $0.413$ unchanged, and effective dimensionality barely moves
(participation ratio $4.879 \to 4.500$ while leakage falls $0.914 \to 0.553$).

It holds on a domain the adversary never saw: Fitzpatrick17k
$0.764 \to 0.391$, while a non-adversarial baseline reaches $0.994$ on the same
images.

\subsection{Models become confidently wrong on unseen domains}
\label{sec:confidently}

At $\lam = 0$, mean maximum-softmax confidence is $0.915$ on ISIC test and
$0.489$ on \texttt{pad\_heldout} --- the model is appropriately hesitant on
unfamiliar data. By $\lam = 0.25$ that gap has vanished ($0.913$ versus
$0.898$), and by $\lam = 2$ it has inverted: the model is \emph{more}
confident on images from a clinic it has never seen ($0.930$) than on its own
held-out test set ($0.910$).

Cross-domain diagnostic accuracy does not improve at any $\lam$. Balanced
accuracy spans $0.249$--$0.291$ against a six-class floor of $0.167$, with no
setting exceeding the $0.291$ obtained at $\lam = 0$ and every setting below
the $0.298$ of the non-adversarial EfficientNet-B3 control; plain accuracy
falls $0.245 \to 0.184$, below the $0.385$ majority-class rate throughout. Macro AUC is roughly preserved ($0.605 \to 0.581$), so
discrimination survives while the decision rule does not: predictions collapse
toward the benign majority class, with nevus predictions rising from 113 to
374 of 716 images and \textbf{actinic keratosis recall falling from $0.244$ to
$0.016$} ($49/\nAK \to 3/\nAK$ images, mean over five seeds). Squamous cell
carcinoma moves the same way ($0.088 \to 0.021$, $6/\nSCC \to 1/\nSCC$) but the
count is too small to carry a claim on its own, and melanoma ($n = \nMEL$)
supports none.

Confidence rising while accuracy falls is what drives out-of-domain
calibration error from $0.247$ to $0.746$. We stress that these are three
views of a single phenomenon rather than three independent results.

\subsection{Leakage does not mediate the collapse}
\label{sec:nonmediation}

Two results rule out the natural hypothesis that removing domain information
is what breaks OOD detection.

First, leakage \textbf{saturates at its chance floor by $\lam = 0.25$}
($0.553$, floor $0.5$) while AUROC continues to fall from $0.508$ to $0.413$
as $\lam$ rises further. The damage continues after invariance is complete.

Second, at the pre-projection backbone, leakage is essentially flat across the
whole sweep ($0.982 \to 0.941$) while backbone Mahalanobis AUROC falls
$0.749 \to 0.475$. Domain information remains linearly decodable and OOD
detection collapses anyway. \textbf{This is the independent evidence that the
five-detector agreement does not provide}, since the backbone is a different
representation.

We therefore report leakage as a measured quantity and never as the
explanatory variable, and plot all results against $\lam$.

\subsection{What the phenomenon is not}

\begin{table}[t]
\centering
\small
\caption{Candidate explanations tested and refuted. Each was a pre-specified
prediction with a defined test.}
\label{tab:refuted}
\begin{tabular}{p{3.2cm}p{5.4cm}p{6.0cm}}
\toprule
Explanation & Test & Result \\
\midrule
Memorisation of the adversarial images & Hold PAD out of the adversary at patient level & \texttt{pad\_heldout} $0.427$ vs \texttt{pad\_adv} $0.407$; unchanged \\
Class composition & Restrict ID to the six shared classes; reweight ID to the OOD class mix & $0.414$ vs $0.409$; reweighting \emph{strengthens} the inversion to $0.240$ \\
Domain specificity & Evaluate on Fitzpatrick17k, never in $\mathcal{L}_{\mathrm{adv}}$ & $0.399$; baseline reaches $0.994$ on the same images \\
Dimensional collapse & Participation ratio and eigenspectrum across $\lam$ & $4.879 \to 4.500$ while leakage $0.914 \to 0.553$; decoupled \\
Latent compression & Latent variance and distance to centroid across $\lam$ & OOD is already $4.5\times$ tighter at $\lam=0$; the adversary \emph{decompresses} it \\
Tail asymmetry & Recompute AUROC after truncating the pooled upper tail & $0.412$--$0.424$ versus $0.427$; medians invert \\
\midrule
Synthetic reproduction & Pre-registered toy model, heavy-tailed ID, grid locked before fitting & No inversion ($\lam{=}0$: $0.693$; $\lam{=}2$: $0.786$) \\
\bottomrule
\end{tabular}
\end{table}

Table~\ref{tab:refuted} summarises six refuted accounts and one failed
controlled synthetic reproduction. We report these prominently because they
are why the remaining claims are credible, and because delimiting the
phenomenon is useful to whoever explains it.

\subsection{A domain monitor is cheap}

A linear domain head on a frozen ImageNet ResNet-50 that has never seen this
data reaches AUROC $0.998$; on the trained baseline, $0.997$
(Table~\ref{tab:monitor}). Detecting the domain is easy from almost any
representation, so retaining a monitor is not in itself a contribution. What
distinguishes the context branch is \emph{concentration}: a single principal
component carries $84.5\%$ of its variance and already separates the domains,
against $0.55$--$0.70$ at $k{=}1$ for every other representation tested.

\begin{table}[t]
\centering
\small
\caption{A domain monitor is obtainable from almost any representation. The
$k{=}1$ column, not the AUROC column, is where the factorized model differs.}
\label{tab:monitor}
\begin{tabular}{lcc}
\toprule
Representation & Linear domain AUROC & AUROC at PCA $k{=}1$ \\
\midrule
Frozen ImageNet ResNet-50 & 0.998 & 0.65 \\
Frozen ImageNet EfficientNet-B3 & 0.994 & 0.55 \\
Trained ResNet-50 baseline & 0.997 & 0.65 \\
Trained EfficientNet-B3 control & 0.991 & 0.70 \\
Context branch $z_c$ & $>0.9999$\,\footnotesize{(2 discordant pairs / 11.6M)} & \textbf{1.00} (84.5\% var.) \\
\bottomrule
\end{tabular}
\end{table}

\subsection{Semantic novelty is not affected in the same way}

With the domain held constant and diagnostic classes held out of training
instead, both branches land between $0.6$ and $0.8$ and the ordering depends
on which class is withheld (DF+VASC: lesion $0.648$, context $0.641$, baseline
$0.687$; SCC: lesion $0.806$, context $0.607$, baseline $0.684$). No near/far
structure organises the difference (Pearson $-0.12$ against a
centroid-similarity index). We report this as a null result. The failure
documented above is specific to covariate shift.

% ------------------------------------------------------------
\section{Discussion}

\subsection{The deployment pathway}

The results compose into a concrete failure sequence. A model is trained with
invariance to acquisition domain, as recommended, and deployed at a new site.
On images from that site the model increasingly predicts the benign majority
class, with actinic keratosis recall falling from $0.244$ to $0.016$
($49$ of $202$ images correctly recalled, then $3$). It does so with
higher confidence than on its own test data. The OOD monitor that should flag
those images for human review scores them as more routine than familiar ones.
And the developer's evaluation shows in-distribution accuracy holding,
in-distribution calibration unchanged, and the leakage metric improved exactly
as intended.

\subsection{The danger zone is where the method succeeds}

Risk scales with \emph{achieved} invariance, not with the nominal weight or
with the presence of an adversarial term. In a separate setting we ran the
same objective with an adversary that failed to train; leakage stayed at
$0.94$ across four orders of magnitude of $\lam$, and the safety failure did
not appear either. A weak adversary produces neither debiasing nor harm. This
matters practically: many deployed configurations may be too weak to induce
invariance, and are correspondingly harmless.

\subsection{Tolerance without generalization}

Full-spectrum OOD detection \cite{yang2023} prescribes that covariate-shifted
in-distribution inputs be accepted rather than flagged. Our models achieve
that prescription and are unsafe because of it. By $\lam = 2$ they accept
out-of-domain images more readily than their own test data, while classifying
them at $0.249$ balanced accuracy against a floor of $0.167$ and reporting
higher confidence than in distribution.

Acceptance and competence are separate properties, and current protocols
evaluate them separately or not at all. We suggest that any claim of tolerance
to covariate shift be accompanied by cross-domain accuracy on the same inputs,
since tolerance of inputs a model cannot classify is a failure mode rather
than a success.

\subsection{Two routes to invariance are not interchangeable}

Wald et al.~\cite{wald2021} showed that enforcing multi-domain calibration
produces invariance. We find the converse fails: enforcing invariance
adversarially destroys out-of-domain calibration, driving ECE from $0.247$ to
$0.746$. The property that one route uses as its foundation is the property
the other route removes.

This suggests a concrete design preference. Where invariance is the goal,
calibration-based objectives may be preferable to adversarial ones, because
they arrive at invariance through a constraint that keeps the out-of-domain
uncertainty signal intact rather than destroying it. We have not tested this
directly and offer it as a hypothesis the present results motivate.

\subsection{Recommendation}

Report OOD detection and out-of-domain calibration \textbf{as a function of
achieved invariance}, not at a single operating point; report cross-domain
accuracy alongside any claim of tolerance to covariate shift; and do not use a
separate dataset as the OOD set without saying that this measures covariate
shift rather than semantic novelty. Practitioners tuning $\lam$ already train
at multiple values, so this is a change in what is reported rather than
additional work. It converts an invisible failure into a visible number.

\subsection{An open mechanism}

Six accounts were tested and refuted and a pre-registered synthetic
reproduction failed. We do not know why the inversion occurs. We note only
that it is a bulk location shift, survives shrinkage estimation, is not
mediated by linear domain-decodability, and appears on domains the adversary
never saw. We report the eliminations so that the search space is smaller for
whoever resolves it.

% ------------------------------------------------------------
\section{Limitations}

\textbf{Absolute cross-domain performance is poor for every model here.}
Plain accuracy on \texttt{pad\_heldout} is $0.245$ for the factorized model and
$0.391$ for the baseline, against a $0.385$ majority rate. No model in this
study is deployable. Our claim concerns the \emph{effect of $\lam$} within a
controlled intervention, not the deployability of any particular model, and
the deployment pathway in the discussion should be read as a mechanism of harm
rather than as a description of a fielded system.

\textbf{One domain pair, one modality contrast.} ISIC is dermoscopy; PAD-UFES
and Fitzpatrick17k are clinical photography. The construct we call domain is
confounded with imaging modality, and Fitzpatrick17k does not resolve this
because it is also clinical photography. Generality beyond dermatology is
untested.

\textbf{Architecture scope.} All results are obtained with a dual-encoder
factorization model. The $\lam$ sweep is a controlled intervention
\emph{within} that architecture; whether single-encoder domain-adversarial
training behaves identically is not established here. An attempt to replicate
on Camelyon17 was \textbf{inconclusive}: the adversary head never trained
(cross-entropy $\ln 3$ from the first epoch at every $\lam$, while a post-hoc
probe read $0.94$), despite a gradient graph verified correct
(head cosine $+1.000$, encoder $-1.000$). \textbf{We report no Camelyon17
result and make no claim that Camelyon17 resists invariance.} We record the
diagnostic --- an adversary failing to train under a verified-correct
reversal --- for future work.

\textbf{Detector non-independence.} The five scoring rules are functions of the
same 16-dimensional latent. The backbone analysis provides the independent
replication.

\textbf{Calibration arithmetic.} Out-of-domain ECE is close to mean confidence
minus accuracy in this setting; the confidence, accuracy and ECE results are
one phenomenon viewed three ways.

\textbf{Seeds.} Five seeds at the endpoints, three at interior points. Effect
sizes are large relative to seed variance, but formal testing across seven
points at $n = 3$ is not claimed.

\textbf{Splits.} ISIC splits are image-level, not patient-level; some patients
may contribute images to both partitions. Because in-distribution accuracy is
compared \emph{across $\lam$ under an identical split}, any inflation applies
equally to every row and does not threaten the reported pattern, though
absolute values may be optimistic. PAD splits are patient-level.

\textbf{Third-domain coverage.} 3{,}887 of 16{,}577 Fitzpatrick17k images were
obtained; the set is web-scraped under mixed copyright, used for evaluation
only, and not redistributed.

% ------------------------------------------------------------
\section{Conclusion}

Domain-adversarial invariance training, applied successfully, leaves a skin
lesion classifier confidently wrong on images from clinics it has never seen
and removes its ability to flag them. In-distribution accuracy, in-distribution
calibration and the domain leakage probe all continue to look correct while
this happens. The model achieves the tolerance of covariate shift that current
OOD taxonomies prescribe, without the generalization that would make that
tolerance safe. We could not explain the inversion after refuting six
candidate accounts, but explaining it is not required to act on it: reporting
OOD detection, out-of-domain calibration and cross-domain accuracy as
functions of achieved invariance costs almost nothing and makes the failure
visible.

% ------------------------------------------------------------
\section*{Data and code availability}
\repolink. ISIC 2019 CC~BY-NC~4.0; PAD-UFES-20 CC~BY~4.0;
Fitzpatrick17k annotations CC~BY-NC-SA~3.0 with images used for evaluation
only and not redistributed.

\section*{Declaration of competing interest}
The authors declare no competing interests.

% ------------------------------------------------------------
\begin{thebibliography}{99}
\bibitem{winkler2019} Winkler, J.K. et al. Association between surgical skin
markings in dermoscopic images and diagnostic performance of a deep learning
convolutional neural network for melanoma recognition. \emph{JAMA Dermatology}
155(10), 2019.
\bibitem{bissoto2020} Bissoto, A., Valle, E., Avila, S. Debiasing skin lesion
datasets and models? Not so fast. \emph{CVPR Workshops}, 2020.
\bibitem{daneshjou2022} Daneshjou, R. et al. Disparities in dermatology AI
performance on a diverse, curated clinical image set. \emph{Science Advances}
8(32), 2022.
\bibitem{ganin2016} Ganin, Y. et al. Domain-adversarial training of neural
networks. \emph{JMLR} 17(59), 2016.
\bibitem{bousmalis2016} Bousmalis, K. et al. Domain separation networks.
\emph{NeurIPS}, 2016.
\bibitem{yang2023} Yang, J. et al. Full-spectrum out-of-distribution
detection. \emph{IJCV}, 2023.
\bibitem{zhang2023} Zhang, J. et al. OpenOOD v1.5: enhanced benchmark for
out-of-distribution detection. \emph{arXiv:2306.09301}, 2023.
\bibitem{ovadia2019} Ovadia, Y. et al. Can you trust your model's
uncertainty? Evaluating predictive uncertainty under dataset shift.
\emph{NeurIPS}, 2019.
\bibitem{wald2021} Wald, Y., Feder, A., Greenfeld, D., Shalit, U. On
calibration and out-of-domain generalization. \emph{NeurIPS}, 2021.
arXiv:2102.10395.
\bibitem{isic2019} Tschandl, P., Rosendahl, C., Kittler, H. The HAM10000
dataset. \emph{Scientific Data} 5, 2018; Combalia, M. et al. BCN20000.
\emph{arXiv:1908.02288}, 2019.
\bibitem{pacheco2020} Pacheco, A.G.C. et al. PAD-UFES-20: a skin lesion
dataset composed of patient data and clinical images collected from
smartphones. \emph{Data in Brief} 32, 2020.
\bibitem{groh2021} Groh, M. et al. Evaluating deep neural networks trained on
clinical images in dermatology with the Fitzpatrick 17k dataset.
\emph{CVPR Workshops}, 2021.
\end{thebibliography}

\end{document}
